\clase{8}{25 de Marzo}{}

Dados $I$ conjunto de índices no vacío, $(S^{I}, \beta^{I})$ y $(S, \beta)$ espacios medibles, sea $J \subseteq I$ finito y no vacío. Luego,

\begin{itemize}
	\item Cilindro finito-dimensional:
	\[ C(J, B) = \{(x_{i})_{i \in I} \in S^{I} \ : \ (x_{j})_{j \in J} \in B \} \quad (B \in \beta). \]

	\item Cilindro especial:
	\[ C \Big( J, \bigtimes_{j \in J} B_{j} \Big) = \bigcap_{j \in J} C(j, B_{j}). \]
	
	\item $\beta^{I} = \sigma(C(i, B) \ : \ i \in I, B \in \beta)$.
\end{itemize}

\begin{remark}
	\begin{enumerate}
		\item Los cilindros especiales forman un $\pi$-sistema (si $I$ es finito, son la clase de rectángulos en $S^{|I|}$).

		\item Los cilintros finito-dimensionales forman un álgebra.

		\item Ambas clases generan a $\beta^{I}$.
	\end{enumerate}
\end{remark}

Dado un PE $(X_{t})_{t \in I}$, vimos que podemos verlo como elemento aleatorio $X$ en $(S^{I}, \beta^{I})$. Luego, definimos la \textit{distribución del proceso} como aquella medida de probabilidad $\mbb{P}_{X}$ inducida por el elemento aleatorio $X$ sobre $(S^{I}, \beta^{I})$. \par
Tener en cuenta que, dado un proceso estocástico con distribución $\mu$ sobre $(S^{I}, \beta^{I})$, siempre podemos construir una versión del proceso en el EdP $(S^{I}, \beta^{I}, \mu)$, tomando como $X$ a la identidad; i.e. $X \colon S^{I} \to S^{I}$ tal que $x \mapsto X(x) = x$. Este $X$ se conoce como la \textit{versión canónica} del PE con distribución $\mu$ (o, simplemente, \textit{proceso canónico}).

\begin{remark}
	Si $S$ es un espacio topológico y $\beta$ es su $\sigma$-álgebra de Borel, también podemos considerar en $S^{I}$ a $\beta(S^{I})$, la $\sigma$-álgebra generada por los abiertos en $S^{I}$ con la \textit{topología producto}.
\end{remark}

\begin{pregunta}
	¿Es cierto que $\beta(S^{I}) = \beta^{I}$?
\end{pregunta}

\noindent \textbf{Respuesta.} En general, NO. Aunque siempre se tiene que $\beta^{I} \subseteq \beta(S^{I})$. Si $\tau(S^{I})$ admite una base numerable e $I$ a lo sumo numerable, entonces SI (por ejemplo, $S^{I} = \R^{\N}$).

\subsection{Construcción de medidas en $S^{I}$}
Sea $I \neq \varnothing$ un conjunto de índices. Definimos
\[ \mbb{F}(I) \coloneq \{J \subseteq I \ : \ J \text{ finito y no vacío}\}.  \]
Dados $J_{1} \subseteq J_{2} \subseteq I$, definimos la proyección $\pi_{J_{2} \to J_{1}} \colon S^{J_{2}} \to S^{J_{1}}$ como
\[ \pi_{J_{2} \to J_{1}}((x_{j})_{j \in J_{2}}) \coloneq (x_{j})_{j \in J_{1}}. \]
En el caso en que $J_{2} = I$, abreviamos $\pi_{J_{1}} \coloneq \pi_{I \to J_{1}}$. \par
Dada una medida $\mu$ sobre $(S^{I}, \beta^{I})$, definimos para $J \in \mbb{F}(I)$ la medida $\mu_{J}$ sobre $(S^{J}, \beta^{J})$ por la fórmula
\[ \mu_{J}(A) \coloneq \pi_{J} \mu(A) \coloneq \mu(\pi_{J}^{-1}(A)) = \mu(\underbrace{C(J,A)}_{"A \times S^{I - J}"}). \]
Llamaremos a la colección $\mbb{F}(\mu) \coloneq \{\mu_{J} \colon J \in \mbb{F}(I)\}$ la \textit{colección de distribuciones finito-dimensionales de} $\mu$. \par
Como ejemplo, notemos que si 
\[ \mu = \mbb{P}_{(U_{n})_{n \in \N}} = \prod_{n \in \N} \mscr{L}|_{[0,1]} \]
entonces,
\[ \mbb{F}(\mu) = \{\mscr{L}|_{[0,1]^{|J|}} \colon J \in \mbb{F}(I)\}. \]

\begin{theorem}[Extensión de Kolmogorov-Daniell]
	Sea $I \neq \varnothing$ un conjunto de índices. Sea $S$ un espacio métrico completo y separable con $\beta$ su $\sigma$-álgebra de Borel. Supongamos que tenemos una familia $\mscr{C} = \{\mu_{J} \colon J \in \mbb{F}(I)\}$ tal que:
	\begin{enumerate}[i)]
		\item para cada $j \in J$, $\mu_{J}$ es una medida de probabilidad en $(S^{J}, \beta^{J})$;

		\item $\mscr{C}$ es \textit{consistente}; i.e. si $J_{1}, J_{2} \in \mbb{F}(I)$ son tal que $J_{1} \subseteq J_{2}$, entonces $\pi_{J_{2} \to J_{1}} \mu_{J_{2}} = \mu_{J_{1}}$ (i.e. $\mu_{J_{2}}(A \times S^{J_{2} - J_{1}}) = \mu_{J_{1}}(A) \quad \forall A \in \beta^{J_{1}}$).
	\end{enumerate}
	Entonces, existe una única medida $\mu$ sobre $(S^{I}, \beta^{I})$ tal que $\pi_{J} \mu = \mu_{J} \quad \forall J \in \mbb{F}(I)$; i.e. $\mscr{C} = \mbb{F}(\mu)$.
\end{theorem}
\begin{proof}[Proof Other Information][Idea]
	Sea $\mcal{A}$ el álgebra de cilindros finito-dimensionales de $S^{I}$. Definimos $\mu \colon \mcal{A} \to [0,1]$ por $\mu(C(J, B)) \coloneq \mu_{J}(B)$. Usando la consistencia, se puede ver que $\mu$ está bien definida. \par
	Luego, si demostramos que $\mu$ es $\sigma$-aditiva en $\mcal{A}$, por Carathéodory se puede extender a todo $\sigma(\mcal{A}) = \beta^{I}$. Esto prueba la existencia. \par
	Por otro lado, como $\mcal{A}$ es un $\pi$-sistema que genera $\beta^{I}$, $\mbb{F}(\mu_{1}) = \mbb{F}(\mu_{2}) \implies \mu_{1},\mu_{2}$ coinciden en $\mcal{A} \implies \mu_{1} = \mu_{2}$ (esto último por Dynkin).
\end{proof}

\begin{corollary}
	Sea $S$ espacio métrico completo y separable, con $\beta$ su $\sigma$-álgebra de Borel. Supongamos que tenemos una familia $\mscr{C} = \{P_{n} \colon n \in \N\}$ tal que
	\begin{enumerate}[i)]
		\item $P_{n}$ es una medida en $(S^{n}, \beta^{n} \coloneq \beta \times \cdots \times \beta)$ para cada $n \in \N$;

		\item $\mscr{C}$ es \textit{consistente}; i.e. $\pi_{\{1,\dots,n+1\} \to \{1,\dots,n\}} P_{n + 1} = P_{n} \quad \forall n \in \N$.
	\end{enumerate}
	Entonces, existe una única medida $\mbb{P}$ en $(S^{\N}, \beta^{\N})$ tal que $\pi_{\{1,\dots,n\}} \mbb{P} = P_{n} \quad \forall n \in \N$.
\end{corollary}
